\section{Lecture 15}

\textbf{Dantzig's Method of Covariance Estimation}
\[\min |\Omega|_1 \st \|\Sigma_n \Omega - I||_\infty \leq \lambda\]
Note that we cannot use the alternate constraint \(\|\Sigma_n - \Omega^{-1}\|_\infty \leq \lambda\) because of the comptuational difficulty of inverting a high-dimensional matrix.

\textbf{Related}: Estimating VAR. \(X_n = A X_{n-1} + \epsilon_n\). How can we estimate \(A\) in the case where \(p > n\)? This is similar to estimating \(A\) in the following sense:
\[X_n X_{n-1}^\intercal = A X_{n-1} X_{n-1}^\intercal + \epsilon_n X_{n-1}^\intercal\]

\[\E[X_n X_{n-1}^\intercal] = A \E[X_{n-1} X_{n-1}^\intercal]\]

Thus,
\[\Gamma = A \Gamma_0\]
where \(\Gamma_\ell\) is the lag \(\ell\) autocovariance matrix. We can then estimate \(A\) via
\[\min |A|_1 \st \|\wh\Gamma_1 - A \wh \Gamma_0\|_\infty \leq \lambda\]
This results in a sparse VAR matrix \(A\). This is useful for applications because it allows us to restrict our scientific interpretation to the smaller number of non-zero coefficients. \citet{WW16} analyze the asymptotic properties of this estimator.

Back to Dantzig: we can split this up into sub-problems. Let \(\omega_j\) denote the \(j\)th column of \(\Omega\). Then the optimization problem splits into
\[\min |\omega_j|_1 \st |\Sigma_n \omega_j - e_j|_\infty \leq \lambda\]
For \(j = 1, \dots, p\). Then we let \(\wh \Omega\) be the symmetrization of \([\wh \omega_1, \dots, \wh \omega_p]\).

\begin{center}\rule{0.5\linewidth}{0.5pt}\end{center}

\textbf{Asymptotic theory for \(\wh \omega_1 := \wh \omega\)}:

\textbf{Step 1:}

Suppose the true value is
\[\omega^* =
\begin{bmatrix}
\omega_{11} \\
\vdots \\
\omega_{p1}
\end{bmatrix}\]
When \(\lambda = 0\), the problem is infeasible (because \(n < p\) and the covariance matrix is rank deficient). When \(\lambda = 1\), the solution \(\Omega \equiv 0\) becomes feasible and the unique solution. In fact, we need to choose \(\lambda\) such that
\[\|\Sigma_n \omega^* - e_1\|_\infty \leq \lambda\]
with high probability. Then
\[\|(\Sigma_n - \Sigma^*)\Omega^*\|_\infty \leq \|\Sigma_n - \Sigma^*\|_\infty \|\Omega^*\|_{L_1}\]
where \(\|\cdot\|_{L_1}\) is the maximum row sum. Suppose \(X_i\) are sub-Gaussian. Then
\[\|\Sigma_n - \Sigma^*\|_\infty \asymp \sqrt{\frac{\log p}{n}}.\]
Thus, we must choose
\[\lambda \gtrsim \Omega\left( \sqrt{\frac{\log p}{n}} \|\Omega^*\|_{L_1} \right)\]
If \(X_i\) is also stationary sub-Gaussian process, then we can get the same result with an extra factor of \(\max_\theta f(\theta)\) where \(f(\theta)\) is the spectral density function. See Basu, Machailidis for details.

\textbf{Step 2}: By the basic inequality,
\[|\wh \omega|_1 \leq |\omega^*|_1\]
By triangle inequality,
\[\begin{gathered}|\Sigma^*(\wh \omega - \omega^*)|_\infty \leq |\wh \Sigma(\wh \omega - \omega^*)|_\infty + |(\wh \Sigma - \Sigma^*)(\wh \omega - \omega^*)|_\infty \\
\leq |\wh \Sigma \wh \omega - e_1|_\infty + |e_1 - \wh \Sigma \omega^*|_\infty + |\wh \Sigma - \Sigma^*|_\infty |\wh \omega - \omega^*|_1\end{gathered}\]
Let \(t = |\wh \Sigma - \Sigma^*|_\infty\) be the maximum entry. We can bound
\[|\wh \Sigma \wh \omega - e_1|_\infty \leq \lambda\]

\[|\wh \Sigma \omega^* - e_1|_\infty \leq \lambda\]

The last term we can bound by \(2\lambda\). So we've proven that \(|\Sigma^*(\wh \omega - \omega^*)|_\infty \leq 4\lambda\). Then we get
\[|\wh \omega - \omega^*|_\infty = |\Omega^* \Sigma^*(\wh \omega - \omega^*)|_\infty \leq 4\lambda |\Omega^*|_{L_1}\]
How can we bound \(|\wh \omega - \omega^*|_1\)? Assume \(0 \leq q < 1\). We need to assume
\[\max_k \sum_{j=1}^p |\omega_{jk}|^q \leq C\]
For instance, when \(q = 0\), this is the assumption that all of the columns are \(C\)-sparse. When \(q \approx 0\), this is a "smallness" condition. Choose a subset
\[A = \{j : |\omega_{j1}| \geq t\}, \qquad B = \{j : |\omega_{j1}| < t\}\]
Then \(|A| \leq C/t^q\).

\textbf{Homework}: Work out the remaining details. Read the CLIME paper \citep{CLM11} if you get stuck.

\begin{center}\rule{0.5\linewidth}{0.5pt}\end{center}

\textbf{High-dimensional regression}: Suppose we have \(Y = W \beta + \epsilon\). We can also do
\[\min |\beta|_1 \st |W^\intercal W \beta - W^\intercal Y|_\infty \leq \lambda\]
The constraint enforces an approximate solution to the normal equations. Using the same approach, we can get
\[|\wh \beta - \beta|_\infty = O_p(t) |\Omega|_1 (1 + |\beta|_1)\]

\begin{center}\rule{0.5\linewidth}{0.5pt}\end{center}

\textbf{Factor models}: this is especially important for high-dimensional time series. Suppose we have data
\[X_{it}, \quad i = 1, \dots, N, t = 1, \dots, T\]
We call \(F_t\) the factor process. Then
\[X_{it} = \lambda_i^\intercal F_t + \epsilon_{it}\]
where \(\lambda_i\) is the \(i\)th loading of the observation \(i\). A special case is
\[X_{it} = F_t + \epsilon_{it}\]
We assume that cross-sectional dependencies are squeezed into \(F_t\), such that \(\epsilon_{it}\) has no cross-sectional dependence. We want to estimate factors and loadings from the data. Assuming normality, the covariance matrix is
\[\one\one^\intercal + I\]
The eigenvalues are \(N + 1\) and \(1, \dots, 1\). We can estimate this with
\[\wh \Sigma_n = \frac{1}{T}\sum_{t=1}^T x_t x_t^\intercal\]
We can compute the eigendecomposition of \(\wh \Sigma_n\) and use a scree plot to find the correct number of factors. We can use Davis--Kahan \citep{DK70} to bound the distance between the eigenvectors. Golub and Van Loan \citep{GVL13} also give a version of this theorem. Also see \citet{YWS15}.
